\documentclass[11pt]{article}
\usepackage[a4paper,margin=1.8cm]{geometry}
\usepackage{fontspec}
\usepackage{hyperref}
\usepackage{enumitem}
\usepackage{titlesec}
\usepackage{setspace}
\usepackage{xcolor}
\usepackage{parskip}

\setmainfont{DejaVu Serif}
\setsansfont{DejaVu Sans}
\setmonofont{DejaVu Sans Mono}
\emergencystretch=2em

\hypersetup{
  colorlinks=true,
  linkcolor=blue!50!black,
  urlcolor=blue!50!black,
  pdftitle={ARI Founder Handbook},
  pdfauthor={GitHub Copilot}
}

\setstretch{1.08}
\setlist[itemize]{leftmargin=1.6em,itemsep=0.25em,topsep=0.35em}
\setlist[enumerate]{leftmargin=1.8em,itemsep=0.25em,topsep=0.35em}
\titleformat{\section}{\large\bfseries}{}{0em}{}
\titleformat{\subsection}{\normalsize\bfseries}{}{0em}{}

\title{ARI Founder Handbook\\\large Foundations, Taste, Execution, and Architecture}
\author{}
\date{March 12, 2026}

\begin{document}

\maketitle
\tableofcontents
\newpage

\section{Purpose of This Handbook}

This handbook is written for a serious builder who wants to train judgment, not just accumulate information.

Its governing idea is simple:

\begin{itemize}
  \item purpose decides what is worth dedicating your life to
  \item will decides whether you can endure difficulty
  \item taste decides what is real, deep, reusable, and worthy
\end{itemize}

Purpose without taste becomes fanaticism. Will without taste becomes waste. Taste is what lets you choose the right mountain.

\section{What Taste Means}

In this context, taste is not preference. It is the ability to distinguish:

\begin{itemize}
  \item signal from noise
  \item mechanism from metaphor
  \item durable leverage from fashionable novelty
  \item trustworthy systems from impressive demos
  \item work that matters from work that merely performs importance
\end{itemize}

The five recurring questions are:

\begin{enumerate}
  \item Does this touch reality, or only benchmark theater?
  \item Does this compound, or reset every hype cycle?
  \item Does this unlock reusable capability, or only a narrow artifact?
  \item Does this deepen our model of mind, biology, society, or infrastructure?
  \item If this succeeds, does it actually matter?
\end{enumerate}

\section{Knowledge System to Rebuild}

If starting from a high-school baseline, the best order is:

\begin{enumerate}
  \item mathematics as the language of structure and change
  \item physics and dynamical systems
  \item computer science as controlled complexity
  \item neuroscience and computation
  \item embodiment and robotics
  \item biology and augmentation
  \item institutions, privacy, and inequality
\end{enumerate}

The unifying lesson is that serious systems are constrained, stateful, adaptive, and embedded in society.

\section{Twelve-Week Operational Plan}

\subsection{How to use the plan}

Each week should include:

\begin{enumerate}
  \item study
  \item implementation
  \item writing
  \item reflection
\end{enumerate}

Each week should also produce one judgment artifact: a short note on what became more real or less convincing.

\subsection{Week 1 to Week 4}

\textbf{Week 1: Functions, variables, and modeling}

Move from school math to causal description. Build tiny simulations of linear growth, exponential decay, and threshold switches. Write one page on why a model is a compressed causal story.

\textbf{Week 2: Derivatives, accumulation, and state}

Learn to think in local change and accumulation through time. Implement Euler integration and observe the effect of step size. Write one page on why discretization already trains engineering taste.

\textbf{Week 3: Vectors, matrices, and noisy observations}

Understand structure and uncertainty together. Implement a matrix-based recurrent update and add sensor noise. Write one page on why real systems are structured and noisy.

\textbf{Week 4: LIF from first principles}

Implement a leaky integrate-and-fire neuron from scratch. Sweep time constant and threshold. Write one page on what LIF preserves and what it discards.

\subsection{Week 5 to Week 8}

\textbf{Week 5: Synapses, spike trains, and coding}

Add synaptic filtering and compare fast and slow decay. Build a two-neuron excitatory or inhibitory motif. Write one page on why not all spikes mean the same thing.

\textbf{Week 6: Recurrent circuits and competition}

Study stabilization, competition, and runaway excitation. Build a small recurrent SNN and identify the difference between strong dynamics and parameter abuse. Write one page on why recurrence is power and danger at once.

\textbf{Week 7: Embodiment and closed loops}

Run the embodied prototype and inspect sensory history, action history, and trajectory. Change one body parameter and observe the shift in behavior. Write one page on why bodies compute too.

\textbf{Week 8: Connectomes and structure-constrained models}

Redraw the prototype network into sensory, interneuron, descending, and motor groups. Justify every connection in words. Write one page on why a wiring diagram is not yet a mind.

\subsection{Week 9 to Week 12}

\textbf{Week 9: Plasticity and local learning}

Implement one toy plasticity rule and track how synaptic weights evolve. Write one page on what information a learning rule actually needs.

\textbf{Week 10: Interfaces and hidden arbitrariness}

Inspect the prototype's sensory mapping and action decoder. Identify which parts are principled and which are hand-tuned. Write one page on why interfaces are where truth leaks.

\textbf{Week 11: Trust, privacy, and human stakes}

Design a toy privacy-first memory system for a personal AI companion. Identify abuse paths. Write one page on why trust is an architecture property before it is a policy statement.

\textbf{Week 12: Synthesis and founder-level clarity}

Write one-page architecture memos for Alive, Realize, and Innocence, plus one page on the shared reusable stack. Write one page on what is worth building.

\subsection{Non-negotiable rules}

\begin{enumerate}
  \item Never let reading replace implementation.
  \item Never let implementation replace writing.
  \item Never let ambition excuse vagueness.
  \item Record confusions explicitly.
  \item Refuse at least one seductive but shallow idea every two weeks.
\end{enumerate}

\section{ARI Taste Standards}

Any major direction should pass six tests:

\begin{enumerate}
  \item reality test
  \item mechanism test
  \item reusability test
  \item trust test
  \item dignity test
  \item durability test
\end{enumerate}

If a direction fails several of these, it should not receive deep commitment.

\subsection{Alive}

\textbf{Human problem}

Loneliness is not only absence of contact. It is absence of witnessed identity, continuity, and trusted presence.

\textbf{Irreducible technical object}

Not a chatbot, but a persistent relational agent with memory, boundaries, continuity, and possibly embodiment.

\textbf{Architecture layers}

\begin{enumerate}
  \item identity layer with stable persona and versioned self-model
  \item memory layer with episodic, preference, and relational memory
  \item interaction layer across text, voice, visual, and possibly embodied channels
  \item user sovereignty layer with real memory editing, deletion, and export
  \item trust and safety layer with crisis-sensitive behavior and clear limits
\end{enumerate}

\textbf{Taste standard}

Prefer continuity over novelty, genuine care over retention tricks, and agency over dependency.

\subsection{Realize}

\textbf{Human problem}

Biological limits appear as disability, decline, suffering, and hard capability ceilings.

\textbf{Irreducible technical object}

Not an enhancement slogan, but a closed-loop augmentation or restoration system grounded in measurable biology.

\textbf{Architecture layers}

\begin{enumerate}
  \item sensing layer for noisy biological signal acquisition
  \item interpretation layer with uncertainty-aware decoding and subject adaptation
  \item intervention layer under strict safety envelopes
  \item clinical and validation layer with evidence tracking and auditability
  \item human factors layer for comfort, adherence, and informed use
\end{enumerate}

\textbf{Taste standard}

Prefer mechanism over spectacle, restoration over theater, and translational seriousness over futurist language.

\subsection{Innocence}

\textbf{Human problem}

Systemic inequality persists when infrastructure is opaque, extractive, and easy to weaponize.

\textbf{Irreducible technical object}

Not a fairness dashboard, but a privacy-first coordination and service infrastructure that preserves agency under power asymmetry.

\textbf{Architecture layers}

\begin{enumerate}
  \item identity and consent layer
  \item data boundary layer with encryption and access control
  \item service orchestration layer across institutions
  \item user interface layer that remains understandable and accessible
  \item governance resilience layer against abuse and capture
\end{enumerate}

\textbf{Taste standard}

Prefer privacy by architecture over policy theater, inclusion as actual usability over rhetoric, and hard trust boundaries over benevolent assumptions.

\section{Shared ARI Core}

The ecosystem should compound through common technical layers:

\begin{enumerate}
  \item identity and continuity
  \item privacy and consent
  \item multimodal interaction
  \item personal memory and retrieval
  \item edge and device integration
  \item safety, auditability, and policy enforcement
  \item adaptation under user control
\end{enumerate}

The strategic principle is simple: products differ at the surface, but trust, memory, and control architecture should compound underneath.

\section{Final Standard}

Choose work that is simultaneously:

\begin{itemize}
  \item real
  \item deep
  \item reusable
  \item worthy
\end{itemize}

Alive should make companionship more real, not more addictive.

Realize should make augmentation more grounded, not more theatrical.

Innocence should make infrastructure more trustworthy, not more rhetorical.

If a technical choice strengthens those standards, it is probably aligned. If it weakens them, it is probably a distraction even if it is fashionable.

\end{document}